%% notes-tex/w037-isobutanol-scrnaseq/sections/5-before-you-grow.tex
%% [[experiments.W037-isobutanol-scrnaseq.strain-selection]]
%% https://github.com/Mjvolk3/torchcell/tree/main/notes-tex/w037-isobutanol-scrnaseq/sections/5-before-you-grow.tex
%%
%% SOURCE: hand-authored, with three exceptions, each traced in place:
%%   - the empty-vector precedent is Fig. 3b and the YZy453 row of
%%     Supplementary Table 1 of the source paper;
%%   - pYZ125 is the pYZ125 row of Supplementary Table 2;
%%   - the one-to-two-copy ddPCR result is Supplementary Table 3.
%% No table in this section. The medium column it refers to is computed in
%% strain_tables.py from whether a strain carries a plasmid.

\section{Before the strains are grown\secstatus{todo}}\label{sec:before}

Three items, in the order they would block work.

\notesec{JC0052 needs an empty vector, and the source paper shows why}
Five of the six selected strains carry a \gene{URA3} CEN plasmid and are held on
SC-ura. \strain{JC0052} carries none, and the CEN.PK2-1C parent is
\gene{ura3-52}, so \strain{JC0052} needs uracil in its medium. Growing all six in
medium containing uracil would remove selection from the other five, and a CEN
plasmid is lost without selection, so by the end of a 48\,h fermentation the
population would carry plasmid-free cells. A single-cell measurement resolves
that as a distinct subpopulation, which is a confound in a run designed to
compare plasmid states.

The source paper met the same problem and solved it by transforming the
plasmid-free strain with an empty vector: Fig.~3b reports YZy91 both with and
without one, and YZy453 is YZy452 carrying empty \texttt{pYZ125}. Doing the same
here means one transformation of \strain{JC0052} with \texttt{pYZ125}, a CEN
\gene{URA3} plasmid carrying P\psub{TDH3} and a multiple cloning site and no
cargo, after which all six run on SC-ura. That strain is not among the twelve, so
the transformation has to happen before the run. Skipping it leaves Axis 1
carrying a medium difference on top of its genetic one, and the two cannot be
separated afterwards.

\notesec{Axis 1 compares a deletion to an overexpression, not to a wild type}
\strain{JC0043} supplies \gene{ILV6} from P\psub{TDH3}, so the pair
\strain{JC0052} against \strain{JC0043} is \gene{ilv6}$\Delta$ against
\gene{ILV6} overexpressed. Neither strain has \gene{ILV6} at its native level. A
comparison against knockout-collection data, where the reference is a wild type
carrying one chromosomal copy under its own promoter, has no matching strain in
this set: the strains with \gene{ILV6} intact, \strain{JC0044} and \strain{JC0045},
sit in a different background and also carry a plasmid copy. Whether that gap
matters is worth settling before the run, because closing it means building an
isogenic \gene{ILV6} wild type and that is a construction rather than a
selection.

\notesec{Transgene copy number is per-isolate on four of the six}
\strain{JC0044}, \strain{JC0045}, \strain{JC0050} and \strain{JC0051} carry
$\delta$-integrated cassettes, so their copy number was fixed when the isolate
was picked and is not specified by the construct. The source paper measured one
to two copies by digital droplet PCR in the strains it checked, which were the
sorted mitochondrial-pathway isolates rather than these four. If expression level
of a pathway gene enters the analysis, these four want their own ddPCR first,
since a twofold difference in copy number is indistinguishable from a twofold
difference in regulation once the reads are counted.
